% SEGMENT OLP-0590-S01
\documentclass[../../../include/open-logic-section]{subfiles}

\begin{document}

\olfileid{sth}{card-arithmetic}{ch}
\olsection{தொடரகக் கருதுகோள்}

முடிவுறாத கண அளவெண்களுக்கும் $2$ இன் அடுக்குகளுக்கும், அதாவது
(\olref[opps]{lem:SizePowerset2Exp} இன்படி) அடுக்குக் கணம்
எடுக்கும் செயலுக்கும், ``நேரடியான'' தொடர்பு உறுதியாக இருந்தால்,
கண அளவெண் அடுக்கேற்றம் \emph{எளிதாக} இருக்கும் என்று முந்தைய
முடிவு (சரியாகவே) சுட்டுகிறது. ஆனால் முடிவுறாத கண அளவெண்கள்
அடுக்குக் கணங்களுடன் அப்படி \emph{செயல்படுகின்றன} என்று
கருத முடியாது.

இதை விளக்கத் தொடங்க, பயனுள்ள சில குறியீடுகளை அறிமுகப்படுத்துவோம்.

\begin{defn}
$\cardsucc{\cardfont{a}}$ என்பது $\cardfont{a}$ ஐவிடக்
கண்டிப்பாகப் பெரிய மீச்சிறிய கண அளவெண் என்க. பின்வரும்
இரண்டு முடிவிலித் தொடர்களை வரையறுக்கிறோம்:
\begin{align*}
	\aleph_{0} &\defis \omega &
	\beth_{0} &\defis \omega\\
	\aleph_{\alpha \ordplus 1} &\defis \cardsucc{(\aleph_{\alpha})} &
	\beth_{\alpha+1} &\defis \cardexpo{2}{\beth_{\alpha}}\\
	\aleph_{\alpha} &\defis \bigcup_{\beta< \alpha} \aleph_{\beta} &
	\beth_{\alpha} &\defis \bigcup_{\beta < \alpha}\beth_{\beta} & \text{$\alpha$ ஓர் எல்லை வரிசையெண் எனில்}.
	\end{align*}
\end{defn}

$\cardsucc{\cardfont{a}}$ இன் வரையறை செல்லுபடியாகும்.
ஏனெனில், ஒவ்வொரு கண அளவெண் $\cardfont{a}$ க்கும்
$\cardfont{a}$ ஐவிடப் பெரிய ஒரு கண அளவெண் இருப்பதாக
\olref[cardinals][classing]{lem:NoLargestCardinal} கூறுகிறது.
% SOURCE CORRECTION TA-STH-036: இங்கு மீச்சிறியதைக் கொடுப்பது வரிசையெண்களின் நன்கு வரிசை.
வரிசையெண்களின் நன்கு வரிசைப்படுத்தல் பண்பு
$\cardfont{a}$ ஐவிடப் பெரியவற்றுள் ஒரு \emph{மீச்சிறிய}
கண அளவெண் இருப்பதை உறுதிசெய்கிறது.
% SOURCE CORRECTION TA-STH-037: மீள்வரையறை கண அளவெண் a க்கு அல்ல, அலெஃப் மற்றும் பெத் தொடர்களுக்கே.
இங்கு $\cardfont{a}$ க்கான தொடர்ச்சி மட்டும் தனியாக
வரையறுக்கப்பட்டது. மேலுள்ள அலெஃப், பெத் தொடர்களின் எஞ்சிய வரையறைகள்
முடிவிலிக்கடந்த மீள்வரையறையால் அளிக்கப்படுகின்றன.

``$\aleph$'' என்ற குறியீட்டை காண்டோர் அறிமுகப்படுத்தினார்.
இது எபிரேய எழுத்துமாலையின் முதல் எழுத்தான \emph{அலெஃப்};
``முடிவிலி'' என்பதற்கான எபிரேயச் சொல்லின் முதல் எழுத்தும் இதுவே.
``$\beth$'' என்ற குறியீட்டை பியர்ஸ் அறிமுகப்படுத்தினார்;
அது எபிரேய எழுத்துமாலையின் இரண்டாம் எழுத்தான \emph{பெத்}.\footnote{1900
டிசம்பரில் காண்டோருக்கு எழுதிய கடிதத்தில் பியர்ஸ் இந்தக் குறியீட்டைப்
பயன்படுத்தினார். ஆனால் $\alpha \geq \omega$ ஆகும்போது
$\beth_\alpha$ இருக்காது என்ற தவறான வாதத்தையும் அதில் அளித்தார்.}
இந்தக் குறியீடுகள் முடிவுறாத கண அளவெண்களைத் தருகின்றன.

% SEGMENT OLP-0590-S02
\begin{prop}
ஒவ்வொரு வரிசையெண் $\alpha$ க்கும் $\aleph_\alpha$,
$\beth_\alpha$ ஆகியவை கண அளவெண்கள்.
\end{prop}

\begin{proof}
இரு கூற்றுகளும் எளிய முடிவிலிக்கடந்த தொகுத்தறிதலால்
கிடைக்கின்றன. \olref[cardinals][classing]{omegaisacardinal}
இன்படி $\aleph_0 = \beth_0 = \omega$ ஒரு கண அளவெண்.
$\aleph_\alpha$ மற்றும் $\beth_\alpha$ ஆகிய இரண்டும் கண
அளவெண்கள் எனக் கொண்டால், $\aleph_{\alpha+1}$ மற்றும்
$\beth_{\alpha+1}$ ஆகியவை வரையறைப்படியே கண அளவெண்கள்.
மேலும் \olref[cardinals][classing]{unioncardinalscardinal}
இன்படி கண அளவெண்களின் கணத்தின் ஒன்றிப்பு ஒரு கண அளவெண்.
\end{proof}
\noindent
மேலும், ஒவ்வொரு முடிவுறாத கண அளவெண்ணும் ஓர் $\aleph$ ஆகும்:

\begin{prop}
$\cardfont{a}$ ஒரு முடிவுறாத கண அளவெண் எனில், ஏதேனும்
தனித்துவமான $\gamma$ க்கு $\cardfont{a} = \aleph_\gamma$.
\end{prop}

\begin{proof}
கண அளவெண்கள் மீது முடிவிலிக்கடந்த தொகுத்தறிதலைப்
பயன்படுத்துவோம். $\cardfont{b} < \cardfont{a}$ எனில்
$\cardfont{b} = \aleph_{\gamma_\cardfont{b}}$ எனக்
கொள்வோம். ஏதேனும் $\cardfont{b}$ க்கு
$\cardfont{a} = \cardsucc{\cardfont{b}}$ எனில்,
$\cardfont{a} = \cardsucc{(\aleph_{\gamma_\cardfont{b}})}=
\aleph_{\gamma_\cardfont{b}+1}$.
$\cardfont{a}$ எந்தக் கண அளவெண்ணின் தொடர்ச்சியும் அல்ல எனில்,
கண அளவெண்கள் வரிசையெண்கள் என்பதால்,
$\cardfont{a} = \bigcup_{\cardfont{b} < \cardfont{a}} \cardfont{b} =
\bigcup_{\cardfont{b} < \cardfont{a}}{\aleph_{\gamma_\cardfont{b}}}$.
ஆகவே $\gamma =
\bigcup_{\cardfont{b} < \cardfont{a}}\gamma_\cardfont{b}$ எனில்,
$\cardfont{a} = \aleph_\gamma$.
\end{proof}

% SEGMENT OLP-0590-S03
ஒவ்வொரு முடிவுறாத கண அளவெண்ணும் ஓர் $\aleph$ என்பதால்,
ஒவ்வொன்றும் ஒரு~$\beth$ ஆகவும் இருக்கிறதா என்று கேட்கலாம்.
அப்படி \emph{இருந்தால்}, முடிவுறாத கண அளவெண்களுக்கும் அடுக்குக்
கணம் எடுக்கும் செயலுக்கும் ``நேரடியான'' தொடர்பு இருக்கும்.
அப்போது பின்வருவது கிடைக்கும்:

\begin{defish}
\emph{பொதுமைப்படுத்தப்பட்ட தொடரகக் கருதுகோள்} (GCH).
எல்லா $\alpha$ க்கும் $\aleph_\alpha  = \beth_\alpha$.
\end{defish}

GCH உண்மையாக இருந்தால், கண அளவெண் அடுக்கேற்றத்தைப் பலவாறு
எளிதாக்கலாம். குறிப்பாக $\cardfont{b} < \cardfont{a}$ எனில்,
$\cardexpo{\cardfont{a}}{\cardfont{b}}$ இன் மதிப்பு
$\cardfont{a}\leq \cardexpo{\cardfont{a}}{\cardfont{b}} \leq
\cardsucc{\cardfont{a}}$ என்ற எல்லைகளுக்குள் இருப்பதாகக்
காட்டலாம். இரண்டில் எது நிகழும் என்பதை, அதாவது
$\cardfont{a} = \cardexpo{\cardfont{a}}{\cardfont{b}}$
அல்லது $\cardexpo{\cardfont{a}}{\cardfont{b}} =
\cardsucc{\cardfont{a}}$ என்பதில் எது என்று,
தீர்மானிக்கும் துல்லியமான நிபந்தனைகளையும் தரலாம்.\footnote{அந்த
நிபந்தனையை \emph{இறுதியிணை} தீர்மானிக்கிறது.}

ஆனால் GCH ஒரு \emph{கருதுகோள்}; \emph{தேற்றம்} அல்ல.
உண்மையில், $\ZFC$ முரண்பாடற்றது எனில்
$\ZFC + \text{GCH}$ உம் முரண்பாடற்றது என்று
\citet{Godel1938} நிறுவினார். ஆனால் $\lnot$GCH ஐயும்
$\ZFC$ உடன் சேர்க்க முடியும் என்பது பின்னர் தெரியவந்தது.
GCH இன் மிக எளிய, முக்கியமான \emph{தனிநிகழ்வை} எடுத்துக்கொள்வோம்:

% SEGMENT OLP-0590-S04
\begin{defish}
\emph{தொடரகக் கருதுகோள்} (CH). $\aleph_1 = \beth_1$.
\end{defish}

$\ZFC$ முரண்பாடற்றது எனில் $\ZFC
+ \lnot\text{CH}$ உம் முரண்பாடற்றது என்று \citet{Cohen1963}
நிறுவினார். ஆகவே தொடரகக் கருதுகோள் $\ZFC$ இலிருந்து
சார்பற்றது.

``தொடரகம்'' என்பது மெய்யெண் கோடு $\Real$ க்கான மற்றொரு பெயர்;
இதனால்தான் இக்கருதுகோளுக்கு அந்தப் பெயர்.
\olref[opps]{continuumis2aleph0} இன்படி
$\card{\Real} = \beth_1$. ஆகவே இயல் எண்களின் கண அளவெண்
$\aleph_0 = \beth_0$ க்கும் தொடரகத்தின் கண அளவெண்
$\beth_1$ க்கும் இடையில் கண அளவெண் இல்லை என்பதே தொடரகக்
கருதுகோள்.

(G)CH ஆனது $\ZFC$ இலிருந்து \emph{சார்பற்றது} என்றால்,
அதன் \emph{உண்மை} பற்றி என்ன சொல்லலாம்? இதைப் பற்றி
\emph{நிறைய} கூறலாம்; கணக் கோட்பாட்டில் அண்மைய பல வளமான
ஆய்வுகள் இச்சிக்கல்களை விசாரிக்கின்றன. ஆனால் இரண்டு சுருக்கமான
கருத்துகள் முக்கியம்.

முதலாவது: இத்தகைய முறையான சார்பின்மை முடிவுகளிலிருந்து GCH
அல்லது CH இன் உண்மை மதிப்பு \emph{நிர்ணயிக்கப்படாதது}
என்று \emph{உடனடியாக} பெற முடியாது. ஒருவேளை நமக்கு இயல்பானதாகத்
தோன்றும் கூடுதல் அடிகோள்களைச் சேர்த்தால், ஏதேனும் ஒரு திசையில்
இக்கேள்விக்குத் தீர்வு கிடைக்கலாம். இதுவே சரியான பதில் என்று
கோடல் தாமே கருதினார்.

இரண்டாவது: CH ஆனது $\ZFC$ இலிருந்து சார்பற்றது என்பது
\emph{வியப்பளிப்பது}; ஆனால் சொல்லின் நேர்பொருளில்
\emph{நம்ப முடியாதது} அல்ல. $\ZFC$ தெரிவிப்பதை மட்டுமே கொண்டு
பார்த்தால், ஒரு கண அளவெண்ணிலிருந்து அதன் தொடர்ச்சிக்குச்
செல்வது, அடுக்குக் கணம் எடுப்பதைவிட நுட்பமான செயலால்
நிகழக்கூடும் என்பதே கருத்து.
% அடுக்குக் கணம் எடுக்கும் செயல் படிநிலைமுறையின் ஒரு படிநிலையிலிருந்து
% அடுத்த படிநிலைக்கு, $V_{\alpha}$ இலிருந்து $V_{\alpha+1}$ க்கு, நகர்த்துகிறது.

இரண்டு கருத்துகளையும் கூறிவிட்டோம். மேலும் அறிய விரும்பினால்,
இவ்விவாதத்தில் ஈடுபட்டுள்ள தத்துவவாதிகளையும் கணிதவியலாளர்களையும்
அணுக வேண்டும்.\footnote{தொடக்கமாக \citet[\S15.6]{Potter2004}
ஐப் படிக்கலாம்.}

\end{document}
